\documentclass[12pt,a4paper,UTF8]{ctexart}
\usepackage[a4paper,top=22mm,bottom=22mm,left=22mm,right=22mm,headheight=15pt]{geometry}
\usepackage{amsmath,amssymb,mathtools,bm}
\usepackage{graphicx,xcolor,booktabs,longtable,array,enumitem,fancyhdr,caption,float,microtype,setspace}
\newcommand{\Needspace}[1]{\par\begingroup\dimen0=#1\relax\dimen2=\pagegoal\advance\dimen2 by-\pagetotal\ifdim\dimen0>\dimen2\newpage\fi\endgroup}
\newcommand{\needspace}[1]{\par\begingroup\skip0=\lastskip\vskip-\skip0\dimen0=#1\relax\advance\dimen0 by\skip0\vskip\dimen0\penalty-100\vskip-\dimen0\vskip\skip0\penalty9999\endgroup}
\usepackage[hidelinks,unicode]{hyperref}
\usepackage[most]{tcolorbox}
\definecolor{ContestBlue}{HTML}{17365D}
\hypersetup{pdftitle={第039套},pdfauthor={竞赛物理教研}}
\setstretch{1.24}
\setlength{\parindent}{2em}
\setlength{\parskip}{0.35em}
\setlength{\emergencystretch}{2em}
\setlist{itemsep=0.35em,topsep=0.4em}
\captionsetup{font=small,labelfont=bf,justification=centering}
\setcounter{secnumdepth}{0}
\pagestyle{fancy}
\fancyhf{}
\fancyhead[L]{\small 第039套}
\fancyfoot[C]{\thepage}
\renewcommand{\headrulewidth}{0.4pt}
\ctexset{
  section={format=\Large\bfseries\color{ContestBlue},beforeskip=1.2em,afterskip=0.7em},
  subsection={format=\large\bfseries\color{ContestBlue},beforeskip=1.0em,afterskip=0.5em},
  subsubsection={format=\normalsize\bfseries,beforeskip=0.8em,afterskip=0.35em}
}
\makeatletter
\newcommand{\tightlist}{\setlength{\itemsep}{0pt}\setlength{\parskip}{0pt}}
\newcommand{\maxwidth}{\ifdim\Gin@nat@width>\linewidth\linewidth\else\Gin@nat@width\fi}
\newcommand{\maxheight}{\ifdim\Gin@nat@height>0.72\textheight 0.72\textheight\else\Gin@nat@height\fi}
\newcommand{\pandocbounded}[1]{\begingroup\setkeys{Gin}{width=\maxwidth,height=\maxheight,keepaspectratio}#1\endgroup}
\makeatother
\begin{document}
\begin{titlepage}
  \centering
  \vspace*{58mm}
  {\Huge\bfseries\color{ContestBlue} 第039套\par}
  \vfill
  \vspace*{28mm}
\end{titlepage}
\section{第 1 题}\label{ux7b2c-1-ux9898}

质量为 \(m\) 的运动者以周期性跳跃方式沿水平地面前进。重力加速度为
\(g\)，脚与地面之间的滑动摩擦系数为 \(\mu\)。人在空中受到线性空气阻力 \[
\boldsymbol f=-k\boldsymbol v.
\]
脚蹬地的时间极短，在这一瞬间可忽略重力冲量。稳定运动时，每次跳跃的轨迹完全相同。求水平方向平均速度的最大值。

\begin{figure}[H]
\centering
\includegraphics[width=0.82\linewidth,height=\textheight,keepaspectratio,alt={题1　周期跳跃与线性空气阻力}]{../figures/question-01.pdf}
\caption{题1　周期跳跃与线性空气阻力}
\end{figure}

\newpage

\section{第 2 题}\label{ux7b2c-2-ux9898}

质量为 \(M\)、边长为 \(a\) 的匀质正三角形薄板静置在光滑水平面上。质量为
\(m\) 的小球以速度 \(v_0\)
沿水平方向撞击薄板一个顶点处外伸的轻质刚性挡板，碰撞完全弹性，挡板尺寸忽略不计。求：

（1）第一次碰撞后小球和三角板的运动状态；

（2）能够发生第二次碰撞时 \(m/M\) 应满足的条件；

（3）第二次碰撞后小球和三角板的运动状态。

\begin{figure}[H]
\centering
\includegraphics[width=0.659\linewidth,height=\textheight,keepaspectratio,alt={题2　小球与可自由运动正三角板的顶点碰撞}]{../figures/question-02.pdf}
\caption{题2　小球与可自由运动正三角板的顶点碰撞}
\end{figure}

\newpage

\section{第 3 题}\label{ux7b2c-3-ux9898}

在 \(N\) 维欧氏空间中，把到某一定点距离不超过 \(r\) 的点集称为半径为
\(r\) 的 \(N\) 维球。已知半径为 \(1\) 的球体积从 \(N=1\) 起依次为 \[
2,\quad
\pi,\quad
\frac{4\pi}{3},\quad
\frac{\pi^2}{2},\quad
\frac{8\pi^2}{15},\quad
\frac{\pi^3}{6},\quad
\frac{16\pi^3}{105},\ldots
\]

（1）写出半径为 \(r\) 的 \(N\) 维球体积；

（2）写出其表面积；

（3）若 \(N\) 维球形液滴的表面能为 \(E_s=\sigma_NS_N\)，求球内外压强差；

（4）若液体在一个方向平移对称，横截面是半径为 \(r\) 的 \((N-1)\)
维球，求柱内外压强差。

\begin{figure}[H]
\centering
\includegraphics[width=0.9\linewidth,height=\textheight,keepaspectratio,alt={题3　单位高维球体积系数的递推}]{../figures/question-03.pdf}
\caption{题3　单位高维球体积系数的递推}
\end{figure}

\newpage

\section{第 4 题}\label{ux7b2c-4-ux9898}

太阳光在六角冰晶内部发生两次折射，在最小偏向角附近形成明亮日晕。冰晶折射率为
\(n=1.309\)。

（1）如图，光线在第一界面的入射角为 \(\alpha\)，求总偏向角
\(\theta(\alpha)\)；

（2）求日晕的圆锥半角；

（3）太阳视角直径为 \(0.5^\circ\)。照片中日晕大圆半径 \(R\)
与太阳光斑小圆半径 \(r\) 之比为 \(15\)。相机严格对焦在人物头部，镜头直径
\(D=6\ \mathrm{cm}\)；日晕弥散可忽略，太阳像存在离焦弥散。求人物头部到相机的距离。

\begin{figure}[H]
\centering
\includegraphics[width=0.72\linewidth,height=\textheight,keepaspectratio,alt={题4　六角冰晶中的两次折射}]{../figures/question-04.pdf}
\caption{题4　六角冰晶中的两次折射}
\end{figure}

\newpage

\section{第 5 题}\label{ux7b2c-5-ux9898}

如图，正四面体的六条棱分别接入电阻 \(R\)、电感 \(L\) 和电容
\(C\)，每一类元件占一组对棱。电源接在顶点 \(1\)、\(2\) 之间，电压为 \[
V(t)=V_0\cos\omega t.
\] 已知 \[
R=\sqrt{\frac LC},
\qquad
\omega=\frac1{\sqrt{LC}}.
\] 求电源输出的总电流。

\begin{figure}[H]
\centering
\includegraphics[width=0.599\linewidth,height=\textheight,keepaspectratio,alt={题5　正四面体六条棱上的交流网络}]{../figures/question-05.pdf}
\caption{题5　正四面体六条棱上的交流网络}
\end{figure}

\newpage

\section{第 6 题}\label{ux7b2c-6-ux9898}

玻尔（Bohr）和惠勒（Wheeler）曾用液滴模型分析原子核裂变。用 \(Z\)
表示核内质子数， \(A\) 表示核子数。把原子核看成球形液滴，其半径为 \[
R=r_0A^{1/3},
\qquad
r_0=1.2\ {\rm fm},
\qquad
1\ {\rm fm}=10^{-15}\ {\rm m}.
\] 核的能量由表面能 \(U_{\rm A}\) 和静电能 \(U_{\rm E}\)
两部分组成，表面能满足 \[
U_{\rm A}=a_sR^2,
\] 其中 \(a_s\) 为待定系数。现有半径为 \(R_0\) 的母核裂变为半径分别为
\(R_1\)、\(R_2\) 的两个子核，两个子核的质量之比为 \[
f=\frac{A_1}{A_2}\geqslant 1,
\qquad
A_1+A_2=A.
\] 母核与子核的质量密度和电荷密度均保持不变。完成下列问题：

（a）估算原子核的质量密度，取质子、中子质量近似相同；

（b）把两个子核的表面能之和写成 \[
U_{\rm A}^{d}=a_sR_0^2\alpha,
\] 求 \(\alpha\) 关于 \(f\) 的表达式；

（c）求母核的静电能 \(U_{\rm E}^{p}\)，以 \(\mathrm{MeV}\) 为单位并用
\(Z\)、\(A\) 表示；

（d）两个子核刚刚分开、仍彼此相切时，求其总静电能 \(U_{\rm E}^{d}\)；

（e）令 \(U_{\rm E}^{d}=\beta U_{\rm E}^{p}\)，求 \(\beta\)；

（f）真正的裂变对应两个子核从相切状态分离到相距无穷远。若该过程中释放的总能量为
\(Q\)，用 \(Q\)、\(Z\)、\(A\)、\(\alpha\) 及只依赖于 \(f\) 的系数
\(\gamma\) 表示 \(a_s\)；

（g）已知 \[
{}^{1}_{0}\mathrm n+{}^{235}_{92}\mathrm U
\longrightarrow
{}^{141}_{56}\mathrm{Ba}+{}^{92}_{36}\mathrm{Kr}
+3\,{}^{1}_{0}\mathrm n
\] 释放 \(Q=173.2\ {\rm MeV}\)，求 \(a_s\)；

（h）把裂变放能条件写成 \[
\frac{Z^2}{A}>C,
\] 求 \(C\)；

（i）求 \(C\) 的最小值，并说明对应质量比 \(f\) 的物理意义。

\begin{figure}[H]
\centering
\includegraphics[width=0.82\linewidth,height=\textheight,keepaspectratio,alt={题6　液滴模型中母核分裂为两个子核}]{../figures/question-06.pdf}
\caption{题6　液滴模型中母核分裂为两个子核}
\end{figure}

\newpage

\section{第 7 题}\label{ux7b2c-7-ux9898}

一根长度为 \(l\)、质量为 \(m\)
的柔软、不可伸长、无弹性的均匀细绳放在光滑水平面上，形成圆心角为
\(\Theta<2\pi\) 的圆弧。沿切线方向向外给绳的一端一个冲量
\(J\)，求极短时间后绳子另一端的速度 \(v_1\)。

为建立模型，依次完成：

（1）三个质量均为 \(m\)
的小球由两根不可伸长轻绳连接，两绳在中间小球处的夹角为
\(\theta\)。以冲量 \(J\)
沿第一段绳的反向作用于第一个小球，求第三个小球的速度 \(v_2\)；

（2）把模型推广到 \(N+1\) 个质量均为 \(m\) 的小球和 \(N\)
段轻绳，求最后一个小球的速度 \(v_3\)；

（3）令离散折线逼近原来的圆弧细绳，求 \(v_1\)。

\begin{figure}[H]
\centering
\includegraphics[width=0.9\linewidth,height=\textheight,keepaspectratio,alt={题7　圆弧细绳的切向冲量传递及离散近似}]{../figures/question-07.pdf}
\caption{题7　圆弧细绳的切向冲量传递及离散近似}
\end{figure}

\newpage

\section{第 8 题}\label{ux7b2c-8-ux9898}

考虑由静质量为 \(m_0\)、彼此几乎无相互作用的质点组成的理想气体，数密度为
\(n\)。气体在动量空间中各向同性。

在非相对论情况下，动量大小落在 \(p\) 到 \(p+dp\) 内的概率为 \[
d\mathcal P
=F_0(p)\,dp
=4\pi p^2
\left(\frac1{2\pi m_0kT}\right)^{3/2}
\exp\left(-\frac{p^2}{2m_0kT}\right)dp.
\]

（1）求温度为 \(T\) 时气体的压强；

（2）考虑弱相对论修正，令 \[
\varepsilon=\frac{kT}{m_0c^2}\ll1
\] 为一阶小量，并设径向动量分布为 \[
F(p)=4\pi p^2A\exp\left(-\frac{E(p)}{kT}\right),
\qquad
E(p)=\sqrt{m_0^2c^4+p^2c^2}.
\] 保留至 \(\varepsilon\) 的一阶，求归一化常数 \(A\)、平均能量
\(\overline E\) 和压强。

积分提示：令 \[
G_j(a)=\int_0^\infty x^j e^{-ax^2}\,dx,
\] 则 \[
\begin{aligned}
G_0&=\frac12\sqrt{\frac\pi a},&
G_1&=\frac1{2a},&
G_2&=\frac14\sqrt{\frac\pi{a^3}},\\
G_3&=\frac1{2a^2},&
G_4&=\frac38\sqrt{\frac\pi{a^5}},&
G_5&=\frac1{a^3},\\
G_6&=\frac{15}{16}\sqrt{\frac\pi{a^7}},&
G_7&=\frac3{a^4},&
G_8&=\frac{105}{32}\sqrt{\frac\pi{a^9}}.
\end{aligned}
\]

\begin{figure}[H]
\centering
\includegraphics[width=0.475\linewidth,height=\textheight,keepaspectratio,alt={题8　各向同性分布在动量空间中的球壳}]{../figures/question-08.pdf}
\caption{题8　各向同性分布在动量空间中的球壳}
\end{figure}
\end{document}
